\section{用感知机构建神经网络}

\subsection{多输出感知机与 Dense 层}

单个感知机只有一个输出。通过将多个感知机并列排放，让所有输入同时连接到每个感知机，就构成了\textbf{Dense 层}（全连接层）。

\begin{figure}[H]
\centering
\includegraphics[width=0.85\textwidth]{slides-images/slide-036.jpg}
\caption{多输出感知机：所有输入密集连接到所有输出，即 Dense 层\protect\footnotemark}
\end{figure}
\footnotetext{Slides 第 36 页。}

对于第 $i$ 个输出神经元：
$$z_i = w_{0,i} + \sum_{j=1}^{m} x_j \cdot w_{j,i}$$

因为所有输入都与所有输出相连，这种层被称为 \textbf{Dense}（密集连接）层或\textbf{全连接层}（Fully Connected Layer）。

在代码中实现 Dense 层非常简洁：

\begin{lstlisting}[caption={TensorFlow 和 PyTorch 中的 Dense 层实现}]
# TensorFlow
import tensorflow as tf
layer = tf.keras.layers.Dense(units=2)

# PyTorch
import torch.nn as nn
layer = nn.Linear(in_features=m, out_features=2)
\end{lstlisting}

\begin{figure}[H]
\centering
\includegraphics[width=0.95\textwidth]{slides-images/slide-037.jpg}
\caption{从零实现 Dense 层：TensorFlow（左）与 PyTorch（右）代码对比\protect\footnotemark}
\end{figure}
\footnotetext{Slides 第 37 页。}

\subsection{单隐藏层神经网络}

在输入层和输出层之间加入一个 Dense 层作为\textbf{隐藏层}（Hidden Layer），就构成了一个单隐藏层神经网络：

\begin{figure}[H]
\centering
\includegraphics[width=0.85\textwidth]{slides-images/slide-039.jpg}
\caption{单隐藏层神经网络结构\protect\footnotemark}
\end{figure}
\footnotetext{Slides 第 39 页。}

隐藏层的计算：
$$z_i = w_{0,i}^{(1)} + \sum_{j=1}^{m} x_j \cdot w_{j,i}^{(1)}$$

输出层的计算（以隐藏层的激活值为输入）：
$$\hat{y}_i = g\left(w_{0,i}^{(2)} + \sum_{j=1}^{d_1} g(z_j) \cdot w_{j,i}^{(2)}\right)$$

其中上标 $(1)$ 和 $(2)$ 分别表示第一层和第二层的参数，$d_1$ 是隐藏层的神经元数量。

\begin{importantbox}{隐藏层的含义}
隐藏层之所以称为"隐藏"，是因为我们在训练数据中只能直接观察到输入和输出，而隐藏层的表示是由网络自动学习得到的，对于外部观察者来说是"不可见"的。隐藏层的每个神经元都学习了输入的某种非线性变换，这些变换共同为最终预测提供了丰富的特征表示。
\end{importantbox}

\subsection{深度神经网络}

通过堆叠更多的隐藏层，就构成了\textbf{深度神经网络}（Deep Neural Network）。"深度学习"中的"深度"正来源于此。

\begin{figure}[H]
\centering
\includegraphics[width=0.85\textwidth]{slides-images/slide-042.jpg}
\caption{深度神经网络：多个隐藏层的堆叠\protect\footnotemark}
\end{figure}
\footnotetext{Slides 第 42 页。}

第 $k$ 层第 $i$ 个神经元的计算公式为：
$$z_{k,i} = w_{0,i}^{(k)} + \sum_{j=1}^{n_{k-1}} g(z_{k-1,j}) \cdot w_{j,i}^{(k)}$$

其中 $n_{k-1}$ 是第 $k-1$ 层的神经元数量。

在代码中构建深度网络同样直观：

\begin{lstlisting}[caption={TensorFlow 和 PyTorch 中的深度网络}]
# TensorFlow
model = tf.keras.Sequential([
    tf.keras.layers.Dense(n1),
    tf.keras.layers.Dense(n2),
    # ...
    tf.keras.layers.Dense(2)
])

# PyTorch
model = nn.Sequential(
    nn.Linear(m, n1), nn.ReLU(),
    nn.Linear(n1, n2), nn.ReLU(),
    # ...
    nn.Linear(nK, 2)
)
\end{lstlisting}

\subsection{本章小结}

从单个感知机出发，通过并列多个感知机构成 Dense 层，再通过堆叠多个 Dense 层构成深度神经网络。每一层的输出作为下一层的输入，逐层提取越来越抽象的特征。深度神经网络的"深度"来源于其多层级结构。

